Harmful memes often express abuse through subtle interactions between images
and overlaid text, and their interpretation may depend on implicit social,
cultural, and event-specific knowledge. Existing methods largely focus on
binary detection, leaving targets, intents, rhetorical strategies, and
supporting evidence implicit. We formulate harmful meme understanding as a
structured interpretation task that jointly predicts harmfulness, target
presence and granularity, primary intent, rhetorical tactics, and image--text
relations. Our five-stage framework extracts global, local, textual, and
cross-modal evidence; acquires external knowledge candidates; verifies their
relevance, validity, and task-specific support; fuses internal evidence with
verified knowledge through task-aware reasoning; and outputs structured
predictions with evidence attribution and a concise rationale. A masked
multi-task objective supports partially available labels, while provenance
records distinguish retrieved documents, generated hypotheses, verified
evidence, and rejected candidates. We train and validate on HarMeme and
evaluate generalization on the held-out Facebook Hateful Memes benchmark under
a leakage-controlled protocol. Beyond harmfulness detection, we separately
assess target, intent, and tactic using schema-validated silver annotations
with field-level coverage. This setup enables a unified evaluation of
cross-domain detection, multi-dimensional interpretation, and auditable
evidence-grounded reasoning.